%% ma: L10-B12-0072
%% lop: 10
%% chuong: 5
%% bai: 12
%% dang: 10.12.2
%% loai: TLN
%% muc: TH
%% dap_an: "5"
%% nguon: "chuyen-de-cac-so-dac-trung-cua-mau-so-lieu-khong-ghep-nhom-toan-10-chuong-trinh-moi.pdf (GV. Phan Nhật Linh) (P01), Bài 01 Phần III Câu 6"
%% kien_thuc: "Lớp 10 Bài 12 (số gần đúng, sai số tuyệt đối, độ chính xác, sai số tương đối, quy tròn số)"
%% trang_thai: chinh-thuc
%% ghi_chu: "Cùng số liệu Bài tập tự luận 3 nhưng khác loại câu; lời dẫn bổ sung cách ước lượng để đáp số duy nhất. Tự duyệt theo yêu cầu giáo viên 22/10/2026."
\begin{ex}
Các nhà toán học cổ đại Trung Quốc đã dùng phân số $\dfrac{22}{7}$ để xấp xỉ số $\pi$. Biết $3{,}1415<\pi<3{,}1416$, sai số tuyệt đối của giá trị gần đúng này nhỏ hơn số có dạng $0{,}00ab$ (với $a$, $b$ là chữ số) tìm được bằng cách ước lượng $\dfrac{22}7<3{,}1429$. Tính $S=a+b$.
\shortans{5}
\loigiai{$\Delta=\dfrac{22}7-\pi<3{,}1429-3{,}1415=0{,}0014$. Vậy $a=1$, $b=4$, $S=5$.}
\end{ex}
